\documentclass[11pt]{article}
% \usepackage[margin=1in]{geometry} % to adjust the page margins
\usepackage[document]{ragged2e}

\usepackage[T1]{fontenc}
\usepackage[round]{natbib}
\usepackage{todonotes}
\usepackage{gensymb}
\usepackage{amssymb}
\usepackage{amsmath}
\usepackage{mathtools}
\usepackage[top=1in,bottom=1in,left=1in,right=1in]{geometry}
\usepackage{alltt}
\usepackage{framed}
\usepackage{xcolor}
\usepackage{comment}
\usepackage{graphicx}
\usepackage{setspace} % For line spacing

\usepackage{xr}[]
\externaldocument{main}
%\usepackage{url}
\usepackage{hyperref}

\renewcommand{\thetable}{\Alph{table}}
\renewcommand{\thefigure}{\Alph{figure}}

\newcommand{\beq}{\begin{equation}}
\newcommand{\eeq}{\end{equation}}
\newcommand{\beqs}{\begin{equation}}
\newcommand{\eeqs}{\end{equation}}
\newcommand{\mx}[1]{\ensuremath{\boldsymbol{\mathrm{{#1}}}}}   % Simple matrix format.


% \newcommand{\review}[1]{\noindent{{\textbf{\textit{#1}}}}\vspace{3mm}}
% \definecolor{colframe}{rgb}{0.01,0.199,0.1}

\newcounter{rtaskno}
\newcommand{\rtask}[1]{\refstepcounter{rtaskno}\label{#1}}

\newcommand{\reviewc}[2]{\begin{leftbar}\noindent\rtask{#1}\ref{#1}. #2\end{leftbar}\vspace*{-4mm}}
\newcommand{\review}[1]{\begin{leftbar}\noindent #1\end{leftbar}}
\newcommand{\answer}[1]{\noindent #1\\[8mm]}
\newcommand{\copyquote}[1]{\vspace{2mm} {\color{blue} #1}\vspace{5mm}}

\newlength{\leftbarwidth}
\setlength{\leftbarwidth}{3pt}
\newlength{\leftbarsep}
\setlength{\leftbarsep}{10pt}
\renewenvironment{leftbar}{%
    \def\FrameCommand{{\color{black}{\vrule width \leftbarwidth\relax\hspace {\leftbarsep}}}}%
    \MakeFramed {\advance \hsize -\width \FrameRestore }%
}{\endMakeFramed}


% \usepackage[colorlinks=True,citecolor=black]{hyperref}

\begin{document}
\onehalfspacing % sets the line spacing to one and a half
\setlength{\parskip}{3mm} % Increase space between paragraphs

\noindent{\bf Poli\v{c}ar and Zupan: 
\\
VERA: Generating Visual Explanations of Two-Dimensional Embeddings \\ via Region Annotation}

\vspace{2mm}
\noindent {\bf COVER LETTER}

\vspace{2mm}


\noindent Dear Editor-in-Chief, Associate Editors,\\
 
\vspace{3mm}

\noindent We would like to submit our manuscript ``\textit{VERA: Generating Visual Explanations of Two-Dimensional Embeddings via Region Annotation}'' to be considered for publication in the ``\textit{IEEE Transactions on Visualization and Computer Graphics}'' as a Regular Paper.

This work presents VERA, a general-purpose method that supports the interpretation of point-based visualizations derived from dimensionality-reduction techniques such as t-SNE, UMAP, and autoencoder-based embeddings. While these approaches enable rich visual representations of complex, high-dimensional data, their interpretation remains challenging. VERA automatically identifies informative regions in two-dimensional embeddings and links them to human-interpretable features, producing concise annotations that summarize embedding structure at a glance. We demonstrate VERA across multiple datasets and report a user study showing that participants using VERA achieved accuracy comparable to manually using an interactive data-mining toolkit while requiring significantly less time and effort, confirming its effectiveness for rapid, static interpretation of complex embeddings.

Our work should be of interest to the visualization, visual analytics, and explainable AI communities, as it bridges automated explanation methods with human-centered visual understanding. Potential reviewers from these fields include Dmitry Kobak (dmitry.kobak@uni-tuebingen.de), Cagatay Turkay (cagatay.turkay@warwick.ac.uk), and Marc Streit (marc.streit@jku.at).

The manuscript has not been published or submitted elsewhere, and both authors approve its submission. An earlier IEEE VIS version was substantially revised following reviewer feedback; their comments and our responses are provided below and in the separate document.

\vspace{3mm}

\noindent Kind regards,
\vspace{3mm}

\noindent Pavlin G. Poli\v{c}ar,
Bla\v{z} Zupan \\ 
\noindent Ljubljana, \today

\newpage

{\bf Revisions Based on Feedback from the IEEE VIS Review} \\[5mm]

A previous version of this manuscript was submitted to the 2024 IEEE VIS Conference. Although the paper was not accepted, we received valuable and detailed feedback, which we used to substantially improve the manuscript that we are now submitting to IEEE TVCG. We would like to thank the reviewers for their constructive comments and suggestions. Below, we provide point-by-point responses to the key comments and concerns raised by the reviewers that remain relevant to this revised version of the paper.

\subsection*{Reviewer 1}

% \reviewc{r1q1}{From my understanding of the authors' system, I disagree with their base assertion that the system does not require some degree of human intervention. To begin the process, the embedding algorithm itself needs to be selected by a user, who must also tune the hyperparameters to produce an embedding that the VERA system can ingest. It can be fairly argued that this step is external to the VERA automated process, but it is undeniable that the system is dependent on receiving an embedding that contains meaningful regions, a fact that is noted by the authors in the final paragraph of Section 3.1.2.}

% \answer{}


\reviewc{r1q2}{The authors state a number of human-tunable parameters that are necessary to adjust to that input embedding to produce a high-quality output from the VERA system, including level of KDE-based contour lines (Section 3.1.2), mapping distortion estimates for each data point (also 3.1.2), and region overlap threshold to determine whether regions should be joined (Section 3.1.3).  It could be argued that each of these parameters has a default and that setting them is optional, but using the defaults does not appear to guarantee a quality, meaningful VERA output.}

\answer{While we agree that including a parameter study would be an interesting addition to the manuscript, given the 12-page limit, we chose to include other sections—such as case studies, a user study, and scalability analysis, which we believe are more informative to readers and provide a stronger validation of our approach. As a general note, we have found our default parameter values to work well on all tested datasets, and we did not adjust the parameters in any of the examples shown in the manuscript.}


\reviewc{r1q3}{While the embedding in Figure 2 that is used throughout the description of the system (Section 3), I found it to be very artificial in both the embedding layout, and the simplicity of the relationships between input features (e.g., note that Country Purchased and Number of Children have identical regions).  Using a more complex data set, such as that provided in the teaser figure and Section 4 would better emphasize the power of the VERA technique (or would further highlight some of the limitations).}

\answer{While we agree that the guiding example used throughout the paper is simple, we designed this dataset specifically so that the different aspects and steps of our procedure could be demonstrated clearly and transparently. Moreover, although simple, this dataset captures many challenges relevant to our method that also arise in real-world data. However, we agree that including additional real-world examples would better demonstrate the capabilities of VERA. Accordingly, we have added several more complex, real-world datasets to case studies section.}


% \reviewc{r1q4}{W4: On the limitations point, the authors do note that the system is limited to datasets in which features have "human-readable, semantic meaning" and also encounters scaling problems as the number of features increases.  I would add several more limitations to these.  (1) the process of discretizing numeric variables into five bins assumes a normal of uniform distribution, whereas other distributions (skewed, bimodal) might be separated in a better way; (2) unless the term "binary indicator variable" is being incorrect used, these five bins still require further breakdown into five separate variables (as do any categorical variables with >2 classes), making the technique effectively limited to binary classification rather than naturally supporting multi-class classification;  (3) the process of merging variables assumes that consecutive order is important (x<5, 5<x<6, x>6), but there are cases such as correlations where values close to -1 and 1 share more in common than those close to 0.}

% \answer{}


\reviewc{r1q5}{The system makes use of both descriptive and contrastive explanations, but only contrastive explanations are clearly called out in the Related Work. Adding the "description explanation" keywords to the appropriate paragraphs in Section 2 will better highlight the differences between VERA and existing systems.}

\answer{We agree with the reviewer and have made this categorization more explicit in the Related Work section by clearly labeling and discussing both descriptive and contrastive explanation approaches.}


%\reviewc{r1q6}{M3:  At the end of Section 3.1.3, it is noted that "Features whose values exhibit no spatial correlation with the embedding space are merged into a single explanatory variable."  What percentage of features fall under this category in a typical embedding? And does this number decrease if users tune the parameters noted in W1 to create a better embedding or in W2 to alter how the embedding is parsed?}

%\answer{}


\reviewc{r1q7}{In Figure 5, the Language plot displays English in two distinct regions of the embedding rather than merging them into a single joined region.  Does this separation have an impact on user understanding, particularly after multiple features are joined to define a region?}

%todo: preveri spodnje in umakni, če se ne strinjaš
\answer{We have clarified this behavior in the manuscript by explaining that VERA intentionally keeps spatially separate occurrences of the same feature value as distinct regions to preserve local structure in the embedding. This clarification has been added to Section 3.2.1 (Descriptive Explanations), where we now note that features may correspond to multiple disjoint regions and that these are handled individually to ensure that region-based annotations remain spatially coherent and interpretable.}

\reviewc{r1q8}{Was there any thought given to extending explanations to describe why regions are (or are not) joined?  It's outside the scope of this system, but I'm thinking of the MNIST Digits dataset in particular. How would a user interpret the decision underlying joining similar numbers like 1 and 7 but not other similar numbers like 6 and 8?}

%todo: preveri tudi tole
\answer{We have clarified in Section 3.1.3 that regions are joined only when they show both high spatial overlap and semantic compatibility of their feature values. This addition explains why visually similar but distinct categories (e.g., digits 1 and 7 vs. 6 and 8) may or may not be merged.}


\reviewc{r1q9}{In the results in Section 5.2, reporting standard deviations along with the means and p-values would be useful context.}

\answer{Agreed. We have now included the corresponding standard deviations and standard errors where appropriate.}


%\subsection*{Reviewer 2}

%todo: spodnje sem zakomentiral, tudi meni ni jasno, kaj je hotel, smo pa tokrat prispevke direktno izpostavili, a ne bi to komentiral

% \reviewc{r2q1}{While the paper overall is well written and motivated, I cannot quite connect the dots for one of the main story points. Understanding embedding projections is challenging, yes. So this paper proposes a way to solve that, but through a static visualization. But the contribution is primarily an algorithm to label regions of a projection.}

% \answer{\textcolor{red}{Blaž, na to ne znam odgovoriti. Introduction sem prestrukturiral in upam, da je zdaj povezava med doprinosom in težavo, ki jo izpostavljamo, bolj razvidna. Verjetno pa lahko ta komentar tudi izpustim iz tega cover letter-ja.}}


\reviewc{r2q2}{The evaluation that was added to the end of the paper after some discussion from the previous reviews then compares the static visualization to an interactive visualization of the same thing, but without the automatic VERA labeling. This evaluation seems odd to me, as my expected comparison would have been how this algorithm compares to other projection labeling algorithms. I totally get the synergy of combining VERA with an interactive tool - it seems like every interactive tool should have the VERA approach built-in and be the default "on load" screen!}

\answer{Thank you for this comment. While we agree that a comparison between VERA and other static labeling algorithms would be valuable, we were unfortunately unable to find any directly comparable approaches available online or in the literature. As noted in the Related Work section, most existing tools are either interactive, limited to a single explanation, or domain-specific, and are therefore not directly comparable to VERA.

After reviewing the literature, we were surprised to find no prior studies comparing static and interactive explanation tools. Therefore, we decided that a comparison between the use of Orange Data Mining toolkit and VERA would be particularly informative for the field. As noted in the Study Design section, many previously proposed interactive approaches can be reproduced in Orange, allowing a fair and representative comparison between VERA and interactive approaches more broadly.}


% \reviewc{r2q3}{Given the length of the algorithm description, I expected the evaluation to be on it compared to other labeling algorithms, either in runtime performance, quality, or measuring user understanding. This relates to my next point below.}

% \answer{As noted in ~\ref{r2q2}, we are unaware of any directly comparable static labelling approaches with which to compare our system.}


\reviewc{r2q4}{Previous reviews ask about how the approaches scales. There has been some new texted added explaining that as features grow to the hundreds of thousands this approach may not work well since the labels will be long. It proposes using LLMs to summarize, which is neat, but my concerns are moreso for how this approach scales not just in features but also in number data points. For example, common use of dimensionality reduction is for ML datasets, which can have millions of points. There is no mention of scale in the current submission, and the only examples we see appear to have just a few hundred points.

Moreover, I have other questions about scalability that isn't touched upon in the paper, which I found odd given it's strong algorithmic contribution. For example, how fast does this take to run? Is it in minutes, seconds, or interactive speed? How was it implemented? Where can it run? I see the code link but understand it's private until acceptance.}

\answer{We thank the reviewer for this useful suggestion. We have now added a dedicated \textit{Scalability and Limitations} section to the manuscript, which includes both a theoretical time-complexity analysis of our approach and empirical benchmarks demonstrating VERA’s runtime and scalability with respect to both the number of samples and features.}


\reviewc{r2q5}{While using dimensionality reduction of larger, labelled, high-dimensional dataset is useful for certain domains (mentioned in the submission), another common application for using this to cluster data for exploration in a unsupervised learning approach. Oftentimes in this case there is no additional metadata/features, or at least none that have interpretable units. It is written off that VERA would need to be extended to tackle this. Is there anything more that could be said here? Any hints, ideas, conjectures of applying VERA to these datasets?}

\answer{As with other general-purpose embedding annotation approaches, VERA relies on a collection of human-interpretable features to generate explanations. In theory, VERA could be extended to other types of data, but domain-specific adaptations would be required—for example, identifying exemplar images for image-based data or representative topics for text-based data. In this manuscript, we focus on a general-purpose approach and leave these extensions for future work, as stated in the conclusion.}


\reviewc{r2q6}{I greatly appreciate the detailed and natural language algorithm description, as it walks a reader through the different types of explanations. Using a synthetic dataset with known features and clusters is a good way to compare results (the figures emphasize this). I will say the amount of new terminology introduced is a little overwhelming at first since it's multiple levels of groups of points (lots of aggregating and merging!). Despite these pros, this section is rather long. Nearly 50\% of the paper is dedicated to this section and description. This balance makes the remainder of the paper feel rushed, and doesn't leave room for more examples, which I explain below.}

\answer{Agreed. We have shortened and streamlined the Methods section while preserving its clear, natural-language description of the approach, and we have added three additional examples to the Usage Scenarios section to better illustrate VERA’s capabilities.}


\reviewc{r2q7}{Since a major contribution of this work is the algorithm and the static labelled projection visualizations, it would be great to see more results on different datasets. This was also highlighted in the previous reviews, and while the new user study includes more than just the synthetic books datasets, we still don't get to see many many visualizations. I'd argue that showing many more examples would strongly help this submission given it's emphasis on static representation.}

\answer{Agreed. We have added VERA explanations for three additional datasets, including the Abalone, Iris, and Penguins datasets, to better demonstrate the method’s versatility and effectiveness.}


\reviewc{r2q8}{It would be nice to have the related work partitioned, say into subsections or using paragraph titles to provide short descriptions for the different topics. While the related work is good, having these would help readability.}

\answer{We have restructured and streamlined the Related Work section to improve readability. Given the current organization and the 12-page limit, adding additional subsection headings would provide limited benefit, so we have decided not to include them.}



%\subsection*{Reviewer 3}

%\reviewc{r3q1}{W1: The introduction could be strengthened by highlighting the techniques in VERA (e.g., contrastive and descriptive explanation selection/ranking) and user study results. While the current introduction provides background and motivation for embedding explanations, it lacks clarity on how VERA addresses mentioned challenges and its effectiveness.}

%\answer{}


% \reviewc{r3q2}{W2: The related work is comprehensive but overlooks recent interactive embedding visualization tools [e.g., 1–3]. For instance, WizMap [1] can automatically generate interactive embedding explanations. How does VERA differ from this work?
% \vspace{3mm}
% \\
% \noindent [1] Wang, Z. J., Hohman, F., \& Chau, D. H. (2023). WizMap: Scalable Interactive Visualization for Exploring Large Machine Learning Embeddings. arXiv preprint arXiv:2306.09328.
% \\
% \noindent [2] Lekschas, F. (2023). Regl-scatterplot: A scalable interactive javascript-based scatter plot library. Journal of Open Source Software, 8(84), 5275.
% \\
% \noindent [3] Sivaraman, V., Wu, Y., \& Perer, A. (2022, March). Emblaze: Illuminating machine learning representations through interactive comparison of embedding spaces. In 27th International Conference on Intelligent User Interfaces (pp.
% 418-432).}

% \answer{Thank you for suggesting the additional references. We have now included a citation to WizMap, which is indeed relevant to our work. We now also mention three approaches that propose methods to compare multiple embeddings in the introductory paragraph of the Related Work section, including a citation to Emblaze. While not directly related to our work, comparing embeddings to one another is an adjacent task to ours. We have decided not to include the citation to Regl-scatterplot, since the publication describes a specific Javascript scatter plot rendering library which is not used in our work.}


\reviewc{r3q3}{VERA assumes that the features in the high dimensional space are interpretable. Many people use t-SNE and UMAP to generate 2D embedding visualizations for neural network's latent representation embeddings. These embeddings have a large number of dimensions (e.g., 1024 dimensions), and each dimension does not map to a human-understandable feature. How does VERA generalize to these ML embeddings? This is briefly discussed in the limitation section. I recommend mentioning the focus on tabular/interpretable features earlier in the paper.}

\answer{Agreed. We have now explicitly stated this limitation in both the Abstract and the Related Work section. In the Abstract, we clarify: ``VERA identifies informative regions in the embedding space and associates them with user-provided human-interpretable features, ...'' and in the Related Work section, we note: ``In domains where features are not directly interpretable -- such as images, compositional data, or deep latent representations -- domain-specific solutions are often required, as in the case of text, images, and drug-prescription data.''}


%\reviewc{r3q4}{W4: Section 3 mentions some configurations in VERA, such as k-means quantization with k=5 and using top-3 explanation selection. How can users adjust these settings in practice? This is not detailed in Sections 4 and 5. It would be great if VERA is more interactive and allows users to easily change configurations through interactions (e.g., slider, drop-down menu).}

%\answer{}


\reviewc{r3q5}{Using Orange data mining toolkit as a baseline model may not be ideal since it lacks embedding explanations. I wonder if it is possible to use related works with embedding explanations as a stronger baseline model. Otherwise, it should be acknowledged in the limitation section.}

\answer{We have now added the following sentence to the Study Design section: ``While not explicitly designed for interpreting embeddings, Orange includes many standard visualization tools for exploring two-dimensional embeddings. Many related approaches based on feature distribution plots and disciminative models can be easily reproduced in Orange using only a handful of widgets.''.

Prior to the study, we provided a short tutorial on how to use Orange in this manner. This is already included in the third paragraph of the Study Design section: ``Before conducting the study, we provided a short tutorial on both Orange and VERA explanations.''}


%\reviewc{r3q6}{M1: What does "satisfying explanations" mean? (Section 3) One might argue that "satisfying explanations" can elicit confirmation bias when interacting with VERA.}

% \answer{}


%\bibliographystyle{abbrvnat}
%\bibliography{main}

\end{document}
